\subsection{Inventario de recursos}

A continuación se mencionan los recursos a los cuales el equipo tuvo total acceso para llevar a cabo la realización de este proyecto:

\subsubsection{Expertos}

\begin{itemize}
  \item Dr. Benjamín Valdés Aguirre
  \item Mtro. Eduardo Daniel Juárez Pineda
  \item Dr. José Antonio Cantoral-Ceballos
  \item Dr. Aurelio Guevara Escobar
  \item Lupita
  \item Colaboradores del CAETEC
\end{itemize}

\subsubsection{Datos}

Para la realización del proyecto y entrenamiento de los modelos, fue indispensable contar con los siguientes datos tabulares en formato .csv extraídos de los dispositivos con los que interactúan las vacas.

Los siguientes datos fueron proporcionados por el socio formador:

\begin{enumerate}
  \item Concentrados de alimento consumido de cada vaca hasta enero de 2026
  \item Producción de leche por sesión de ordeño de cada vaca hasta enero de 2026
  \item Registros detallados de los ordeños de cada vaca, patadas e inventario total de marzo a junio de 2025
\end{enumerate}

También se recolectaron varias imágenes de distintas vacas para resolver el problema del BCS. Estos datos fueron recopilados por distintos integrantes de los equipos del proyecto de la concentración.

\subsubsection{Recursos de cómputo (hardware)}

\noindent
\begin{tabularx}{\textwidth}{>{\raggedright\arraybackslash}X>{\raggedright\arraybackslash}X>{\raggedright\arraybackslash}X>{\raggedright\arraybackslash}X}
  \toprule
  Recurso & Especificaciones & Propietario / Acceso & Uso en el proyecto \\
  \midrule
  Computadora de escritorio & CPU: Intel Core i7-12700k\newline GPU: NVIDIA GeForce RTX 4070 Ti 12GB\newline RAM: 32GB\newline 2TB SSD NVMe M.2 & Emiliano Camacho Ponce & Entrenamiento de los modelos \\
  \addlinespace
  Cámara de teléfono & Modelo: iPhone 14\newline Cámara: 12 MP con campo de visión de 120 grados.\newline Formato: 4:3 & Alfredo Alejandro Soto Herrera & Teléfono con el cual se tomaron las imágenes \\
  \bottomrule
\end{tabularx}

\medskip

Cabe mencionar que para la recopilación total de imágenes, se utilizaron otros dispositivos móviles de otros miembros de distintos equipos de la concentración. Se siguió un estándar donde todos los dispositivos utilizados fueran iPhone.

\subsubsection{Recursos de software}

\noindent
\begin{tabularx}{\textwidth}{>{\raggedright\arraybackslash}X>{\raggedright\arraybackslash}X>{\raggedright\arraybackslash}X>{\raggedright\arraybackslash}X}
  \toprule
  Categoría & Herramienta & Versión & Uso en el proyecto \\
  \midrule
  Lenguaje & Python & x.x & Desarrollo general \\
  Deep Learning & TensorFlow & x.x & Diseño y entrenamiento de las CNNs \\
  Análisis de datos & NumPy, pandas, Matplotlib & x.x & Exploración, métricas y visualización \\
  Entornos de desarrollo & Jupyter, VS Code & - & Experimentación \\
  Control de versiones y gestión & Git/GitHub & - & Trabajo en equipo y seguimiento de tareas \\
  \bottomrule
\end{tabularx}
